\section{Conclusion}

This thesis investigated whether the advanced legal reasoning capabilities of a large, reinforcement-learning-aligned teacher model can be transferred into lightweight student Small Language Models (SLMs) suitable for secure, on-premise deployment. To this end, we designed and empirically validated a three-phase distillation and alignment framework, comprising Offline Logit Distillation, On-Policy Knowledge Distillation, and Diagnosis-Driven Preference Optimization, applied to Qwen3 students at the 0.6 billion and 1.7 billion parameter scales on Vietnamese legal Natural Language Inference (NLI), Multiple-Choice (MCQ), and Syllogism tasks.

The principal findings of this research can be summarized as follows:
\begin{enumerate}
    \item \textbf{A complete three-phase pipeline was realized and validated.} The sequential combination of offline distillation \parencite{hinton2015distillingknowledgeneuralnetwork}, on-policy distillation \parencite{agarwal2024onpolicydistillationlanguagemodels}, and DPO alignment \parencite{rafailov2024directpreferenceoptimizationlanguage} progressively transferred, refined, and aligned the student's reasoning. The strongest in-distribution configuration, the 1.7B model after DPO, attained an Overall score of 0.5280 on the internal Split-Test, with an NLI Accuracy of 0.8156 and a Syllogism ROUGE-L of 0.5684.
    \item \textbf{Quantized adaptation is essential at small scale.} On the 0.6B student, Quantized Low-Rank Adaptation (QLoRA) acted as an implicit regularizer that prevented catastrophic forgetting: the hard-label cross-entropy loss decreased steadily (0.5 $\rightarrow$ 0.25) under QLoRA, in contrast to the upward drift (0.3 $\rightarrow$ 0.7) observed under standard LoRA, and raised the Syllogism ROUGE-L by roughly 33\%.
    \item \textbf{On-policy distillation mitigates exposure bias.} By training the 1.7B student on its own rollouts under teacher guidance, the on-policy phase improved NLI Accuracy from 0.7305 to 0.7870 in-distribution and yielded the highest NLI Accuracy (0.920) and Syllogism ROUGE-L (0.372) among all distilled models on the out-of-distribution VLSP2025 Public-Test.
    \item \textbf{LLM-based diagnosis produces cleaner alignment data.} Replacing rule-based error diagnosis with an LLM-as-a-Judge protocol increased the rate of genuinely acceptable rollouts from 11.1\% to 20.7\%, demoted spurious ``lucky guesses'' (reducing the \texttt{RISKY} pool from 34.5\% to 15.5\%), and improved detection of substantively flawed reasoning (from 40.5\% to 63.5\% \texttt{WRONG}), directly enhancing the quality of the DPO preference set.
    \item \textbf{Model scale is the decisive factor.} The 1.7B student consistently outperformed its 0.6B counterpart across every phase, frequently by a factor of $1.5$–$2\times$. Crucially, the benefits of on-policy distillation and DPO alignment were realized only at the 1.7B scale; at 0.6B these same techniques destabilized the model, indicating that a minimum capacity threshold is required for self-generated trajectories and preference penalties to be productive.
\end{enumerate}

Collectively, these results demonstrate that a carefully staged distillation pipeline, anchored by quantized adaptation and semantically-grounded preference diagnosis, can endow a compact 1.7B legal SLM with substantial reasoning competence. At the same time, the study surfaces an important caveat regarding the trade-off between in-distribution specialization and out-of-distribution generalization, which we examine in the following section.
